\chapter{Designing Over-the-Air Updates}
\label{ch:ota}\index{OTA}\index{firmware update}

The production chapter (Chapter~\ref{ch:production}) listed over-the-air update as
the largest of the four gaps and left it at that. This chapter is the design it
deferred: a concrete sketch of how OTA would be built on \emph{this} boot path --
the partition layout, the one genuinely hard mechanism, the boot-time slot
selection, and the rollback state machine. It is a design, not a description: none
of this is implemented today. But it is grounded in what the project actually has --
a readable Ada bootloader, a working HTTPS transport, and pure-Ada signing
primitives -- so it is a buildable plan rather than a wish.

The shape of the problem is worth stating up front because it explains every
decision that follows. An update has to replace the running firmware \emph{with the
running firmware doing the replacing}. You cannot overwrite the image you are
executing, so there must be two slots and a way to choose between them at boot; and
because the thing doing the writing is fallible -- the power can drop mid-write, the
download can be truncated, the new image can be subtly broken -- every step must be
recoverable. OTA is less a feature than a careful exercise in never being caught
without a working image.

\section{What already exists}

The \emph{transport} is done. The W5500 network stack (Chapter~\ref{ch:networking})
and the pure-Ada TLS client (Chapter~\ref{ch:tls}) can already fetch bytes over
authenticated HTTPS -- the weather example does exactly this against a real server.
An OTA client needs no new networking; it needs somewhere to put the bytes and a
bootloader that will choose them.

The \emph{signing} primitives are done too. The TLS work implements RSA-PSS and
ECDSA-P256 verification in pure Ada; image authentication (\S\ref{ota:trust}) reuses
them rather than inventing anything.

What is missing is everything between ``bytes arrive'' and ``the new image runs'':
the slots, the ability to write the boot flash at all, and the bootloader logic to
select and validate.

\section{Two slots and a selector}\index{partition table}

Today the flash carries a single \emph{factory} application partition at
\texttt{16\#1\_0000\#}. OTA replaces it with the standard A/B arrangement:

\begin{center}
\begin{tabular}{l l}
\toprule
\textbf{Partition} & \textbf{Role} \\ \midrule
\texttt{ota\_0}   & application slot A \\
\texttt{ota\_1}   & application slot B \\
\texttt{otadata}  & two-copy selector: which slot, and its state \\
\texttt{factory}  & optional immutable fallback image \\
\bottomrule
\end{tabular}
\end{center}

At any moment one slot holds the running image and the other is the spare that the
next update is written into. The \texttt{otadata} region is the small but critical
piece: it records \emph{which} slot to boot and \emph{what state} it is in. It is
kept in two copies with a monotonically increasing sequence number and a checksum,
so that a power loss while updating the selector itself leaves at least one valid
copy -- the bootloader takes the highest-numbered copy that checksums, and never
trusts a half-written one.

\section{The hard part: writing the flash you boot from}\index{flash!self-write}

The mechanism everything hinges on is the one the project has so far avoided
entirely: erasing and programming the ESP32-S3's own boot flash from the running
application. It is hard for a specific reason. The CPU executes its code from that
same flash, in place, through the instruction cache (Chapter~\ref{ch:perf}); while
the MSPI controller is busy erasing or programming, it cannot also serve a cache
miss, so \emph{any} instruction or constant fetched from flash mid-operation would
hang or fault.

The discipline that makes it safe is therefore exacting, and it is the same one
ESP-IDF's flash driver uses:

\begin{itemize}
\item the erase/program routine, and everything it calls, must run from
  \textbf{IRAM} -- not flash -- so no fetch touches the busy MSPI;
\item the \textbf{instruction and data caches are disabled} for the duration, then
  invalidated and re-enabled after;
\item on this \textbf{dual-core} part, the \emph{other} core must be parked in an
  IRAM spin-loop first, or it will fetch from flash the instant the cache goes down;
\item interrupts are masked because an ISR resident in flash would do the same.
\end{itemize}

The actual erase and program are ROM services -- the mask ROM exposes sector-erase
and page-program entry points alongside the \texttt{esp\_rom\_spiflash\_read} the
bootloader already uses -- so this is not a from-scratch SPI-NOR driver. It is a
\emph{critical-section wrapper} in Ada and a little assembly: park the sibling core,
disable caches, call the ROM, restore. That wrapper is small, but it is the
brick-risky heart of the feature -- a stray write to the bootloader region corrupts
boot itself (recoverable only over USB) -- which is why it deserves to be built and
tested in isolation, against a spare sector, before any of the rest is wired up.

\section{Choosing a slot at boot}

The bootloader gains a step it does not have today. Instead of loading a fixed
address, \texttt{boot\_main} reads \texttt{otadata}, picks the selected slot, and --
this is the part the current loader skips entirely -- \emph{validates the image
before jumping to it}: the magic byte as now, but also a hash over the image, and
(if secure boot is in play) a signature (\S\ref{ota:trust}). Validation is what lets
the loader \emph{fall back}: a slot that fails its hash is not booted; the loader
tries the other slot, or the immutable \texttt{factory} image, rather than jumping
into rubble.

One state in \texttt{otadata} drives rollback: \emph{pending-verify}. When the
application flips the selector to a freshly written slot, it marks that slot
pending-verify rather than confirmed. The bootloader will boot a pending-verify slot
exactly once; if the new firmware comes up healthy it calls back to mark the slot
\emph{confirmed}, and the pending flag clears. If instead the new image crashes or
hangs before confirming -- and a watchdog (Chapter~\ref{ch:production}) is what
guarantees ``hangs'' becomes ``resets'' -- the next boot sees a still-pending slot,
concludes the update failed, and reverts to the previous good slot. The device can
brick its \emph{application} and still recover on the next power cycle.

\section{The update, end to end}

Put together, the application side is a small state machine over the pieces above:

\begin{enumerate}
\item \textbf{Fetch} the new image over HTTPS into the inactive slot, writing it
  through the IRAM flash-program critical section as it streams (no need to buffer
  the whole image -- a sector at a time, with internal SRAM as the DMA staging
  buffer, since PSRAM cannot be a DMA endpoint, Chapter~\ref{ch:perf}).
\item \textbf{Verify} the written slot in place: re-hash it from flash and check the
  signature, so a truncated or tampered download is caught \emph{before} it can be
  selected.
\item \textbf{Arm}: write the selector to the new slot, marked pending-verify, and
  reboot.
\item \textbf{Confirm}: the new image, once it has proven itself (network back up, a
  self-test passed), marks its slot confirmed. If it never does, the next boot
  rolls back.
\end{enumerate}

An anti-rollback counter in eFuse closes the obvious attack on this loop -- pushing
an \emph{older}, validly-signed image with a known flaw -- by letting the bootloader
refuse any image whose version is below a monotonic floor.

\section{Trust, and why OTA cannot stand alone}\label{ota:trust}\index{secure boot}

The one thing OTA must never do is broaden the attack surface it runs over. A loader
that downloads code and jumps to it is, without verification, a remote-code-execution
mechanism wearing a helpful face. So image authentication is not an optional
companion to OTA -- it is part of it. At a minimum every image carries a SHA-256 the
bootloader checks; for anything connected, that hash is signed, the bootloader
verifies the signature against a public key fused into the device, and the matching
secure-boot eFuse makes the mask ROM enforce the same on the bootloader itself
(Chapter~\ref{ch:production}). This is exactly where the pure-Ada ECDSA and RSA-PSS
from the TLS chapter earn a second use: the verification a browser does on a
certificate is the verification a bootloader does on a firmware image, the same code
serving both.

\section{What it would cost}

The honest accounting: the transport and the cryptography are in hand, but the
middle is real work, in rough order of difficulty. The flash-self-write critical
section is small but delicate and brick-risky, and is the right first spike --
nothing else can be tested until a sector can be erased and re-read safely. The
bootloader changes (slot selection, validation, the pending-verify state machine)
are moderate and must be got exactly right because a bug there is a bug in the one
component with no fallback. The partition and \texttt{otadata} format is
straightforward once the model above is fixed. And the application state machine is
ordinary networking code over the new flash primitive. It is a multi-phase feature
that wants secure boot built alongside it rather than after -- but every phase is
buildable from parts this project already contains, which is the whole reason it is
worth sketching rather than importing.
